\section{Symbols and Assumptions}

\subsection{Symbols}

The notation used throughout the paper is fixed in Table~\ref{tab:symbols}: rows follow
the order of first appearance, one symbol means exactly one quantity, and the unit is
written in parentheses inside the meaning cell.

\begin{longtable}{@{}>{\centering\arraybackslash}p{0.16\textwidth}>{\raggedright\arraybackslash}p{0.80\textwidth}@{}}
  \caption{Symbols.}
  \label{tab:symbols} \\
  \toprule
  Symbol & Meaning \\
  \midrule
  \endfirsthead
  \caption{Symbols (continued).} \\
  \toprule
  Symbol & Meaning \\
  \midrule
  \endhead
  \bottomrule
  \endfoot
  $x_1$ & [quantity] (unit)\\
  $x_2$ & [quantity] (unit)\\
  $y(t)$ & [quantity] (unit)\\
  $\theta$ & [quantity] (rad)\\
  $\Delta t$ & sampling step (s)\\
  $I$ & [objective / index] (unit)\\
\end{longtable}

\subsection{Assumptions}
\begin{enumerate}
  \item [Simplification assumption shared by all questions.]
  \item [Parameter / modelling assumption shared by all questions.]
  \item \textbf{(Problem 2)} [What the scenario fixes: knowns, decision variable.]
  \item \textbf{(Problem 3)} [What is unmeasurable, what replaces it.]
\end{enumerate}

% 【Section writing prompt】
% 位置：Question Analysis 之后、Model 之前——读者必须先见到每个符号，才遇到用到它的公式。
% 符号表：两列，单位写在含义里；一个符号一个含义；循环下标不进表；数值放参数表。
% 假设：场景相关的那几条要标出服务哪一问（如上例 3、4），套话式假设（"本文假设…合理"）全部删。
